\section{Tripod kernels and integral Soul\'e generators}
\label{sec:tripod-integral}

\subsection{Correct field formulation and the odd regular theorem}

For a prime $p$, write
\[
 X=\mathbb P^1_{\mathbb Q}\setminus\{0,1,\infty\},\qquad
 \rho_p:G_{\mathbb Q}\longrightarrow
 \operatorname{Out}\bigl(\pi_1^{\mathrm{et}}(X_{\overline{\mathbb Q}})^{(p)}\bigr),
\]
\[
 F=\mathbb Q(\mu_p),\quad K=\mathbb Q(\mu_{p^\infty}),\quad
 \Omega_p^*=\overline{\mathbb Q}^{\ker\rho_p},
\]
and let $\Omega_p$ be the maximal pro-$p$ extension of $F$ unramified outside $p$. Thus the desired equality is
\[
 \Omega_p^*=\Omega_p,
 \quad\text{equivalently}\quad
 \ker\rho_p=\operatorname{Gal}(\overline{\mathbb Q}/\Omega_p).
\]
The kernel is a subgroup of $G_{\mathbb Q}$; it is not itself the maximal ramification Galois group. At $p=2$ the ramification convention at infinity is specified in Theorem~\ref{thm:pro2-kernel}.

Put $G_p=\operatorname{Gal}(\Omega_p^*/K)$, with the Ihara weight filtration $F^mG_p$, and
\[
 \mathfrak g_p=\bigoplus_{m\geq1}F^mG_p/F^{m+1}G_p.
\]
For odd $m\geq3$, choose $\sigma_m\in F^mG_p$ whose Soul\'e value generates the actual image ideal $\kappa_m(F^mG_p)$. Let
\[
 \psi:\operatorname{Lie}_{\mathbb Z_p}(s_3,s_5,\ldots)
       \longrightarrow\mathfrak g_p,
 \qquad s_m\longmapsto\sigma_m,
\]
and let $\Psi$ be the corresponding homomorphism from the free pro-$p$ group on these generators.

\textbf{Sharifi's exact theorem.} For an \emph{odd regular} prime $p$, $\Psi$ is onto. If rational Deligne--Ihara holds, then $\psi$ and $\Psi$ are isomorphisms and $\Omega_p^*=\Omega_p$ \cite[Theorem~1.1, author's preprint p.~2]{Sharifi2002}. Brown supplies the rational conjecture \cite{Brown2012,Brown2019}. Therefore both the kernel-field equality and integral odd free generation hold for odd regular primes. This resolves the correctly formulated known row LM6~A.44 without step~17. Ishii's recent primary account restates exactly this scope \cite[Theorem~1.1 and Remark~1.1.1, p.~2]{IshiiKernel2026}.

\subsection{A separate deduction for the prime \texorpdfstring{$2$}{2}}
\label{sec:pro2-proof}

The following result is a deduction developed and checked in this run. We did not locate it as a stated theorem in the literature searched, and we make no novelty claim. Its proof uses established arithmetic and rational Deligne--Ihara results together with the complete argument below. It does not use the proposed step--17 estimate. The initiating researcher, two other agents, and the parent researcher each read the whole argument; the check ledger distinguishes those reviews from the cited literature inputs and from formal verification.

\begin{theorem}[K17-P2-MOMENT]\label{thm:pro2-kernel}
Let $X=\mathbb P^1_{\mathbb Q}\setminus\{0,1,\infty\}$ and $\Pi=\pi_1^{\mathrm{et}}(X_{\overline{\mathbb Q}})^{(2)}$. Let $\Omega_2$ be the maximal pro-$2$ extension of $\mathbb Q$ unramified at every finite prime other than $2$, with complexification at the real place allowed. If $\Omega_2^*$ is the fixed field of the kernel of the outer representation
\[
 \rho_2:G_{\mathbb Q}\longrightarrow\operatorname{Out}(\Pi),
\]
then $\Omega_2^*=\Omega_2$.
\end{theorem}

The convention at infinity is part of the statement. The theorem concerns faithfulness of a pro-$2$ action; it does not assert that its integral weight Lie algebra is freely generated by odd Soul\'e elements.

\paragraph{Factorization through the maximal ramification group.}
The image of $\rho_2$ is pro-$2$. Here is a direct justification that does not invoke an odd-prime formulation. The geometric group $\Pi$ is free pro-$2$ of rank two, and its Frattini quotient is $\mathbb F_2(1)^2$. The Galois action on this quotient is trivial. For any finite $2$-group $Q$, the kernel of
\[
 \operatorname{Aut}(Q)\longrightarrow
 \operatorname{Aut}(Q/\Phi(Q))
\]
is a $2$-group. To see this directly, fix a basis of $Q/\Phi(Q)$ of size $d$. Its lifts form a set of cardinality $|\Phi(Q)|^d$, and every such tuple generates $Q$ by the defining property of the Frattini subgroup. The displayed kernel acts freely on this set, since an automorphism fixing a generating tuple is the identity. Its order therefore divides a power of $2$. Use the characteristic finite quotients obtained from the descending Frattini series of $\Pi$. They are cofinal among its finite quotients, and an automorphism trivial on $\Pi/\Phi(\Pi)$ induces an automorphism trivial on the Frattini quotient of each of them. Consequently the corresponding kernel in $\operatorname{Aut}(\Pi)$ is pro-$2$. Inner automorphisms lie in that kernel, so the kernel in $\operatorname{Out}(\Pi)$ is also pro-$2$.

At every odd prime, the marked curve $(\mathbb P^1,\{0,1,\infty\})$ has good reduction with disjoint marked sections. Prime-to-residue-characteristic specialization therefore makes inertia act trivially on its outer pro-$2$ fundamental group. Equivalently, one may use the all-primes factorization recalled in Hain--Matsumoto, \S8.4, p.~27 \cite{HainMatsumoto2001}. Combining unramifiedness outside $2$ with the pro-$2$ image gives an induced continuous representation of
\[
 P=\operatorname{Gal}(\Omega_2/\mathbb Q).
\]
It suffices to prove that representation faithful.

\paragraph{The arithmetic source.}
Gras--Jaulent's Theorem/Definition~0(iv) and Theorem~2 give
\[
 P\simeq C_2*\mathbb Z_2
\]
for the unrestricted-at-infinity group \cite[pp.~1--3]{GrasJaulent2008}. The specialization of their $2$-rationality hypothesis to $\mathbb Q$ is elementary. A rational number whose valuations at all odd primes are even is, modulo rational squares, $\epsilon2^e$ with $\epsilon\in\{1,-1\}$ and $e\in\{0,1\}$. Being a square in $\mathbb Q_2$ forces $e=0$ by valuation and $\epsilon=1$ because $-1$ is not a $2$-adic square. The global and local $2$-primary root-of-unity groups both have $2$-rank one. Thus their criterion applies.

Kronecker--Weber identifies the full abelianization with the cyclotomic character
\[
 \chi:P\longrightarrow\mathbb Z_2^\times.
\]
Write $P=\langle c\rangle*\langle\gamma\rangle$, where $c^2=1$ and $\langle\gamma\rangle\simeq\mathbb Z_2$. Then $\chi(c)=-1$. Replacing $\gamma$ by $c\gamma$ if necessary is an automorphism of the free product and ensures that
\[
 u=\chi(\gamma)\in1+4\mathbb Z_2.
\]
Since $\chi$ is the abelianization, $u$ is a topological generator of $1+4\mathbb Z_2$.

The subgroup $P_i=\ker(\chi\bmod4)$, corresponding to $\mathbb Q(i)$, is free pro-$2$ on $\gamma$ and $c\gamma c$. One can verify this without an infinite-rank subgroup theorem: let the involution exchange the two generators $a,b$ of a free pro-$2$ group. Its semidirect product with $C_2$ has the universal property of $C_2*\mathbb Z_2$, with $a=\gamma$ and $b=c\gamma c$. Its index-two kernel is the asserted free group. Set
\[
 h=(c\gamma c)\gamma^{-1},\qquad H=\ker\chi.
\]
Then $(\gamma,h)$ is another free pro-$2$ basis of $P_i$, and $H$ is the kernel of the map $P_i\to\mathbb Z_2$ sending $\gamma$ to $1$ and $h$ to $0$. For $x\in\mathbb Z_2$, write
\[
 h_x=\gamma^xh\gamma^{-x}.
\]

\paragraph{Finite quotients and measure-valued Magnus coefficients.}
Let $N_n$ be the closed normal subgroup of $P_i$ generated by $\gamma^{2^n}$. Then
\[
 P_i/N_n\simeq C_{2^n}*\mathbb Z_2.
\]
The kernel $H_n$ of its map to $C_{2^n}$ is free pro-$2$ on the $2^n$ elements $h_j$, $j\bmod2^n$. Indeed, the free pro-$2$ group on those symbols, with the cyclic permutation action, has semidirect product with $C_{2^n}$ possessing precisely the universal property of $C_{2^n}*\mathbb Z_2$.

The subgroups $N_n$ decrease to $1$. In any finite continuous $2$-group quotient of $P_i$, the image of $\gamma$ has finite $2$-power order, so $N_n$ lies in its kernel for all sufficiently large $n$. Compactness also makes the map to the inverse limit surjective: compatible quotient cosets form a nested family of nonempty closed subsets of $P_i$. Hence
\[
 P_i\simeq\varprojlim_nP_i/N_n,
 \qquad
 H\simeq\varprojlim_nH_n.                                  \tag{P2.1}
\]
The transition maps fold the finite bases by reduction of $j$ modulo $2^n$. The elements $h_x$ topologically generate $H$. To check this last point explicitly, the kernels of $H\to H_n$ decrease to $1$, and compactness implies that every open normal subgroup contains one of them. Every finite quotient of $H$ therefore factors through an $H_n$, whose indicated generators are images of the $h_x$.

The finite-rank integral Magnus maps are injective. A primary locator is \cite[\S5.1, p.~21, (5.1)]{MinacPasiniQuadrelliTan2020}, whose mod-$2$ statement implies the integral injection. There is also a short direct verification. The augmentation ideal of $\mathbb F_2[Q]$ for a finite $2$-group $Q$ is nilpotent: choose a central element $z$ of order two, use $(z-1)^2=0$, and induct on $Q/\langle z\rangle$. Thus any map from a finite-rank free pro-$2$ group to $Q$ is recovered by substituting $X_j\mapsto\overline g_j-1$ into the mod-$2$ Magnus series. An element with trivial expansion maps to $1$ in every finite quotient and is therefore $1$.

The degree-$r$ coefficients of the compatible integral expansions in (P2.1) belong to
\[
 M_r=\varprojlim_n\mathbb Z_2[(\mathbb Z/2^n)^r]
     =\mathbb Z_2[[\mathbb Z_2^r]],
\]
the bounded $\mathbb Z_2$-valued measures on $\mathbb Z_2^r$. With multiplication given by external product and concatenation, they define an injective homomorphism
\[
 \theta:H\longrightarrow1+\prod_{r\geq1}M_r,
 \qquad \theta(h_x)=1+\delta_x.                             \tag{P2.2}
\]
Injectivity follows from the finite-rank maps and (P2.1). Consequently every nontrivial element has a first finite nonzero Magnus degree. Completion in the $(2,I)$ topology creates no element with all homogeneous coefficients zero: the completed coefficient algebra is the displayed product.

\paragraph{The actual rational image.}
Let $B$ be the proalgebraic closure over $\mathbb Q_2$ of the rational outer image induced by the tripod representation, and let $U$ be its prounipotent radical. Hain--Matsumoto, \S8.3 and Proposition~8.3, Theorem~8.4, pp.~25--26, gives an extension
\[
 1\longrightarrow U\longrightarrow B\longrightarrow\mathbb G_m\longrightarrow1
\]
whose reductive quotient is induced by $\chi$, and a continuous map
\[
 \rho_U:H\longrightarrow U(\mathbb Q_2)                    \tag{P2.3}
\]
with Zariski-dense image \cite{HainMatsumoto2001}. The rational action is induced functorially from the pro-$2$ action; a trivial pro-$2$ outer automorphism gives a trivial rational outer automorphism. Their Appendix~A.6, A.9 and A.10 records the relevant comparison. We have used the actual arithmetic image, not an unrelated motivic group.

Brown's rational Deligne--Ihara theorem and Hain--Matsumoto's graded comparison give, with positive arithmetic weights,
\[
 \operatorname{gr}\operatorname{Lie}U
 \simeq\operatorname{Lie}_{\mathbb Q_2}(s_3,s_5,s_7,\ldots),
 \qquad\operatorname{wt}(s_m)=m.                           \tag{P2.4}
\]
The exact outer-action statement for every prime is Brown \cite[\S2.1, Conjecture~1, p.~2; \S2.3, Theorem~2.1 and following paragraph, p.~4]{Brown2019}; the original theorem is \cite[Theorem~1.1 and its consequence, p.~950]{Brown2012}.

We spell out the completion step. In a separated complete positively filtered Lie algebra whose associated graded is free with only finitely many generators below each weight, choose lifts of the free graded generators. The resulting map from the weight-completed free Lie algebra induces an isomorphism on associated gradeds. Induction through successive finite weight quotients proves it is an isomorphism on each quotient; inverse limits give an isomorphism of the completed Lie algebras. A Levi splitting of the negatively weighted extension, as in \cite[Lemma~3.7, p.~9]{HainMatsumoto2001}, permits weight-homogeneous choices. Thus the completed universal enveloping algebra of $\operatorname{Lie}U$ is the weight-completed free associative algebra on the $s_m$. Its finite weight quotients are finite-dimensional.

Let $\ell_m:U(\mathbb Q_2)\to\mathbb Q_2(m)$ be the $m$th abelian coordinate, and put $a_m=\ell_m(\rho_U(h))$. Conjugation by $\gamma$ acts on that coordinate by $u^m$, because the unipotent part of a lift acts trivially on abelianization. Hence
\[
 \ell_m(\rho_U(h_x))=a_mu^{mx}.                             \tag{P2.5}
\]
Every $a_m$ is nonzero. Otherwise that coordinate would vanish on all the topological generators $h_x$, hence on $H$, contradicting the density in (P2.3) and the nonzero $m$th abelian quotient in (P2.4). No integral primitive normalization is asserted here.

\paragraph{Substitution at a finite weight cutoff.}
Suppose $g\neq1$ lies in the kernel of (P2.3). Let $r$ be its first nonzero Magnus degree, and let
\[
 0\neq\mu=\mu_r(g)\in\mathbb Z_2[[\mathbb Z_2^r]].
\]
In the completed free associative algebra put $Y_x=\rho_U(h_x)-1$. Its constant term is zero and its part of generator length one is
\[
 \sum_{m\geq3,\ m\text{ odd}}a_mu^{mx}s_m.
\]
Here generator length is distinct from arithmetic weight and from geometric depth.

Substitution $\delta_x\mapsto Y_x$ is well-defined and continuous coefficient by coefficient. At target weight at most $N$, only Magnus degrees at most $\lfloor N/3\rfloor$ contribute. The remaining coefficient functions on compact powers of $\mathbb Z_2$ are continuous and bounded in $\mathbb Q_2$. Integration against a bounded $\mathbb Z_2$-measure is defined by uniform approximation by locally constant functions, after scaling into $\mathbb Z_2$ if necessary. The same approximation proves continuity, multiplicativity and Fubini's identity. This construction needs no lattice uniform across all weights. The substitution and $\rho_U$ agree on every $h_x$ and are continuous homomorphisms on each finite weight quotient, so topological generation makes them agree on $H$. No universal property of a free pro-$2$ group is applied to a noncompact target.

The coefficient of the ordered word $s_{m_1}\cdots s_{m_r}$ in the image of $g-1$ is therefore
\[
 \left(\prod_{j=1}^ra_{m_j}\right)
 \int_{\mathbb Z_2^r}\prod_{j=1}^ru^{m_jx_j}\,d\mu(x).
\]
Lower Magnus degrees are zero, and higher ones have generator length greater than $r$. Since $\rho_U(g)=1$ and every $a_m\neq0$, we obtain
\[
 \int_{\mathbb Z_2^r}\prod_{j=1}^ru^{m_jx_j}\,d\mu(x)=0
 \quad(m_j\geq3\text{ odd}).                              \tag{P2.6}
\]
The variables remain ordered; this is not a symmetrized moment.

\paragraph{Analytic separation.}
The Amice transform
\[
 A_\mu(T_1,\ldots,T_r)
   =\int_{\mathbb Z_2^r}\prod_{j=1}^r(1+T_j)^{x_j}\,d\mu(x)
\]
lies in $\mathbb Z_2[[T_1,\ldots,T_r]]$ and converges on the open unit polydisc. It is injective as a transform of bounded measures: its coefficients are the moments against products of binomial polynomials, the multivariate Mahler basis. Equivalently it is the canonical isomorphism of the multivariable Iwasawa algebra with the power-series algebra.

For odd $m\geq3$ put $t_m=u^m-1$. These points lie in $4\mathbb Z_2$, and the distinct points $t_{1+2^j}$ converge to $u-1$. Distinctness follows from the infinite order of $u$. Equation (P2.6) says that $A_\mu$ vanishes on the full Cartesian grid of these points. The one-variable nonarchimedean identity theorem applies on the closed disc of radius $1/2$, contained in the radius-one convergence disc. Fixing grid values in all but one variable and applying that theorem successively in each variable shows $A_\mu=0$. Therefore $\mu=0$, a contradiction.

It follows that $\rho_U$ is injective on $H$. The kernel of the pro-$2$ tripod action is contained in $H$ because geometric abelianization records $\chi$. Its elements are also trivial under the induced rational outer action, so they are $1$. This proves Theorem~\ref{thm:pro2-kernel}.

\paragraph{What this proof does and does not decide.}
The proof uses the explicit arithmetic presentation $C_2*\mathbb Z_2$, not merely abstract freeness of $H$. It settles the field equality with the stated convention at infinity. It does not prove integral odd generation of the weight Lie algebra, and it does not transfer to irregular primes without a substitute for its linear-moment separation. The latter obstruction is made explicit below.


\subsection{Four meanings of an integral generator}
\label{sec:four-integral-meanings}

The following claims differ and must be kept separate:
\begin{enumerate}
\item[(A)] A nonzero integral odd class $\sigma_m\in\mathfrak g_{p,m}$ exists.
\item[(B)] The full-domain character
\[
 \kappa_m:G_{\mathbb Q(\mu_{p^\infty})}^{\mathrm{ab}}
       \longrightarrow\mathbb Z_p(m)
\]
is onto.
\item[(C)] An element of the weight subgroup $F^mG_p$ has $\kappa_m$-value $1$ in a fixed integral normalization.
\item[(D)] The map $\psi$ above is an isomorphism of graded $\mathbb Z_p$-Lie algebras.
\end{enumerate}

For odd $p$, (A) is already supplied by Soul\'e nonvanishing and Sharifi's integral construction \cite[Proposition~2.3, pp.~5--6]{Sharifi2002}. It is primitive relative to the \emph{actual image ideal}, which need not be $\mathbb Z_p$. At $p=2$, existence also follows from Brown and the actual lattice: multiply a nonzero rational odd generator by a sufficiently large power of $2$ to put it in $\mathfrak g_{2,m}$, then choose a representative in $F^mG_2$. This establishes existence, independently of the kernel theorem and of step~17.

For odd $p$, all of (B), for odd $m\geq3$, hold if and only if Vandiver holds at $p$. More precisely, each character is nonzero, and it is onto if $m\equiv1\pmod{p-1}$; simultaneous surjectivity is equivalent to Vandiver \cite[Theorem~1.1(1)--(3), p.~2]{IshiiKummer2024}. The theorem is attributed there to the classical Soul\'e and Ichimura--Sakaguchi work. We verified this precise primary-paper locator; we did not verify an original Ichimura--Sakaguchi theorem number and do not invent one. On all of $G_{\mathbb Q}$ the class is a cocycle, while on the cyclotomic kernel it is an ordinary continuous homomorphism.

There is a decisive restriction issue in (C). Sharifi's integral stable-derivation comparison gives, on $F^mG_p$,
\[
 \kappa_m=(p^{m-1}-1)(m-1)!\lambda_m^{(p)},
 \qquad \lambda_m^{(p)}\text{ integral}
 \quad\cite[(2.10), p.~7]{Sharifi2002}.
\]
Thus
\[
 \kappa_m(F^mG_p)\subseteq(m-1)!\mathbb Z_p(m).              \tag{I.1}
\]
At $p=3,m=5$, this image is contained in $3\mathbb Z_3$. Yet $3$ is regular, so (B) and (D) hold. Consequently full-domain surjectivity, weight-subgroup primitivity, and integral generation cannot be conflated. Formula (I.1) does not dispute the cohomological index formula for the full-domain cyclotomic class; it concerns a further restriction.

For (D), odd regular primes are settled positively by Sharifi--Brown. The prime $37$ gives a negative answer despite Vandiver and the corresponding vanishing of all the $K_{4k}$-groups at that prime. The proof follows.

\subsection{A counterexample to integral odd generation at \texorpdfstring{$37$}{37}}
\label{sec:integral37}

Let $\mathfrak h_p=\operatorname{im}\psi$. McCallum--Sharifi's notation in \S9 agrees with this choice: their ``Deligne conjecture in degree $i$'' means injectivity of $\operatorname{gr}^i\psi$ \cite[p.~24]{McCallumSharifi2003}. Brown gives precisely this injectivity, because the rational map is injective and every graded piece of the integral free source is torsion-free.

\begin{proposition}[T-INTEGRAL37]\label{prop:37-counterexample}
At $p=37$, Vandiver holds, all $K_{4k}(\mathbb Z)\otimes\mathbb Z_{37}$ vanish for $k\geq1$, and all full-domain odd Soul\'e characters are onto. Nevertheless
\[
 (\mathfrak g_{37}^{\mathrm{ab}})_{32}\neq0,
 \qquad
 \mathfrak h_{37,32}\subsetneq\mathfrak g_{37,32}.
\]
In particular $\mathfrak g_{37}$ cannot be generated over $\mathbb Z_{37}$ by homogeneous elements all of odd weight, and is not the free graded integral Lie algebra on one generator in every odd weight at least three.
\end{proposition}

\begin{proof}
The only irregular index for $37$ is $32$. Vandiver at $37$ can be verified by the small exact modular certificate accompanying this report, using the criterion of Buhler--Harvey \cite[\S4, pp.~4--5]{BuhlerHarvey2011}: take $q=149$ and $\zeta=107$, whose square root in the order-$37$ subgroup is $16$. The half-exponent normalization is corrected in Hart--Harvey--Ong, \S7, p.~13; the detailed certificate above uses exactly that correction. The required product is $V=48\pmod{149}$ and $V^{(149-1)/37}=142\not\equiv1\pmod{149}$. The computation uses exact integer arithmetic and was independently cross-checked; it is not a formal proof certificate. The modern K-theoretic comparison recorded elsewhere in this report gives the asserted $K_{4k}$ vanishing. The full-domain character assertion follows from \cite[Theorem~1.1(3)]{IshiiKummer2024}.

McCallum--Sharifi's Theorem~7.5 proves that their pairing at index $32$ for $37$ is nonzero. Their Corollary~8.3, p.~23, makes a coefficient explicit: $e_{3,32}$ is nonzero, represented by $-11$ up to a nonzero normalization scalar. Their Theorem~9.7, p.~28, applies when Vandiver holds, $m$ is the least irregular even index, rational Deligne holds through degree $m$, and some $e_{i,m}\neq0$ has $i<m/2$. All hypotheses hold with $(p,m,i)=(37,32,3)$, by Brown and $3<16$.

For clarity, the proof of that cited theorem identifies the mechanism. In degrees below $32$, the integral map is bijective. Thus the weight-$32$ commutator part of $\mathfrak g$ equals that of $\mathfrak h$. The nonzero cup coefficient, through their Proposition~9.6, prevents bijectivity at $32$; Brown's injectivity makes the defect a failure of surjectivity. Their theorem consequently gives both displayed conclusions.

If a graded Lie algebra is generated by homogeneous odd-weight elements, its abelianization is generated by their odd-weight images and has no even-weight component. The nonzero degree-$32$ abelianization contradicts this. This proves the proposition.
\end{proof}

The authors' Magma computation in Theorem~7.5 is used as a published theorem; we did not rerun it. The deduction above, including every hypothesis of Theorem~9.7, was read in full. It does not say that $\mathfrak h_{37}$ is nonfree: rational injectivity makes it a free integral subalgebra. It says that its degree-$32$ lattice is proper. The graded pieces of $\mathfrak g_{37}$ remain torsion-free; their Lie abelianization can have $37$-torsion.

Most importantly, Proposition~\ref{prop:37-counterexample} does not refute $\Omega_{37}^*=\Omega_{37}$. An arithmetic source and its faithful image can have these same extra integral directions. This is a counterexample to the proposed integral-generation implication, not a counterexample to the kernel-field equality.

Sharifi also identifies the first possible defect under Vandiver: if $\Psi$ is not onto and $m$ is its first defective weight, then $m$ must be even and $p\mid B_m$ \cite[Theorem~4.1, p.~10]{Sharifi2002}. The odd-weight alternative would violate Vandiver. This explains why removing the plus class-group obstruction does not eliminate all integral-generation obstructions.

\subsection{Irregular primes: the sharper remaining gate}
\label{sec:irregular-gate}

Write $H_p=\operatorname{Gal}(\Omega_p/K)$ and $M_p=H_p^{\mathrm{ab}}$, with the conjugation action of $\Gamma=\operatorname{Gal}(K/F)\simeq\mathbb Z_p$ and $\Lambda=\mathbb Z_p[[\Gamma]]$. A useful correction is that $H_p$ is already free pro-$p$ even at irregular primes. Iwasawa's original Theorem~2 states that the maximal $S$-ramified pro-$p$ group over a cyclotomic $\mathbb Z_p$-field is free exactly when its $\mu$-invariant is zero \cite[pp.~273--274]{Iwasawa1981}. Ferrero--Washington applies to $K/\mathbb Q$, which is abelian. This does not make the finite-base group $\operatorname{Gal}(\Omega_p/F)$ free.

Let $A_\infty$ be the direct limit, under extension of ideals, of the $p$-parts of the class groups in the cyclotomic tower. Kummer duality gives
\[
 M_p^+\simeq(A_\infty^-)^\vee(1)
 \quad\cite[Corollary~3.3.14 and Proposition~3.4.7]{SharifiIwasawaNotes}.
                                                               \tag{I.2}
\]
The finite-level minus class groups inject into $A_\infty^-$, by the same source's Proposition~3.4.1. At an irregular prime that minus class group is nonzero, so $M_p^+\neq0$. With $\mu=0$, this is a nonzero finite-rank free $\mathbb Z_p$-module but a torsion $\Lambda$-module. The distinction between $p$-torsion and $\Lambda$-torsion is essential.

Every odd linear Soul\'e coordinate kills $M_p^+$: complex conjugation acts by $+1$ on the source and $-1$ on the target $\mathbb Q_p(m)$, so equivariance gives twice the value equal to zero. Thus the degree-one separation used in the pro-$2$ proof fails on a concrete nonzero module. A lift of a nonzero class in $M_p^+$ can have a nonzero first source Magnus coefficient while all target generator-length-one coordinates vanish. Its image may still be nontrivial at higher length. Abstract freeness of $H_p$, weak Leopoldt, or $\mu=0$ alone therefore cannot repeat that proof.

A sharper continuation is to control the full nonlinear lifts of the $\Gamma$-relations defining this torsion module. Their first nonlinear terms are cup products; the degree-$32$ result at $37$ detects one such direction through a bracket. Kernel equality requires detection of every remaining direction and its iterated commutators, or an equivalent separation theorem for the full noncommutative relation ideal. A class-group order or characteristic ideal alone does not specify those nonlinear terms. The present run supplies no such all-weight theorem, so the irregular-prime field equality remains unresolved here. Neither Proposition~\ref{prop:37-counterexample} nor (I.2) gives a negative answer.

\subsection{The integral comparison at \texorpdfstring{$2$}{2}}

Theorem~\ref{thm:pro2-kernel} and rational Deligne--Ihara do not prove surjectivity of
\[
 \operatorname{Lie}_{\mathbb Z_2}(s_3,s_5,\ldots)
      \longrightarrow\mathfrak g_2.
\]
Nonzero integral odd classes exist as explained above. Choosing odd classes whose rational images form a free generating system gives an injective integral free source. Equality with the actual lattice requires an additional argument. Theorem~\ref{thm:2-allweights} below supplies that argument in every weight: it constructs primitive adapted odd generators, proves integral freeness, and uses the exact outer-filtration comparison. It does not prescribe unit values for a separately normalized Soul\'e family.

There is a useful rational normal form, but it does not remove that question. Choose a Levi subgroup containing the image of $c$. Since $chc=h^{-1}$, the element $a=\log h$ has only odd total weights. Writing $\gamma=vu$, the element $\exp(a/2)v$ is $c$-fixed. A further even unipotent conjugation can remove this factor because $u^w-1\neq0$ in $\mathbb Q_2$ for every positive weight $w$, giving a rational form $\gamma=\exp(-a/2)u$. Division by $2$ and inversion of $u^w-1$ need not preserve an integral lattice. This explains the denominators in the finite-weight lattice calculation. The subsequent parity induction uses actual group commutators to remove them where needed and supplies the integral free basis.

\subsection{Mixed Tate motives and the current \texorpdfstring{$\mathfrak{grt}_1$}{grt1} statement}

Deligne--Goncharov's category over $\mathbb Z$ has rational coefficients. In the attached LaTeX, \S1.6, Propositions~1.7 and~1.9, and Proposition~2.3 establish the rational Ext and free-Lie statements; \S2.4 records the noncanonical nature of free lifts \cite{DeligneGoncharov2005}. The base $\operatorname{Spec}\mathbb Z$ specifies ramification conditions. It does not replace $\mathbb Q$-coefficients by $\mathbb Z$ or $\mathbb Z_p$. Rational $\operatorname{Ext}^2$ vanishing erases torsion and is not an integral vanishing result.

The rational Grothendieck--Teichm\"uller status must also be updated from current primary sources. Hirose's October2026 manuscript, arXiv:2610.23719v1, Theorem~4 and \S1.2, gives a Hoffman spanning result over $\mathbb Z_{(2)}$ for the confluence quotient and derives rational equality of the de Rham mixed-Tate motivic Galois group with GRT \cite{Hirose2026}. We checked the statement and its stated implication, not the whole proof. This is a separate current rational result. It does not identify the arithmetic $\mathbb Z_p$-image lattice, and it does not invalidate the counterexample at $37$. An integral associator lattice, a $\mathbb Z_{(2)}$-span of motivic periods, and the lattice of an arithmetic pro-$p$ Galois image are different objects unless a comparison theorem is supplied.

The supplied round-two war-game file mentions B3 and S9 but does not reproduce their original round-one rows. We have therefore used their subject as specified in the prompt and have not assigned an invented exact locator to missing text.

% Bibliography keys needed in this fragment and pro2_section.tex:
% Sharifi2002: Relationships between conjectures on the structure of pro-p Galois groups unramified outside p; https://www.math.ucla.edu/~sharifi/publ/galstr.pdf
% Brown2012: Mixed Tate motives over Z, Annals175(2012)949-976; https://annals.math.princeton.edu/wp-content/uploads/annals-v175-n2-p10-p.pdf
% Brown2019: From the Deligne-Ihara conjecture to multiple modular values; https://arxiv.org/pdf/1904.00179
% HainMatsumoto2001: Weighted completion of Galois groups and Galois actions on the fundamental group of P1-{0,1,infinity}; https://arxiv.org/pdf/math/0006158
% GrasJaulent2008: Note on 2-rational fields; https://arxiv.org/pdf/0805.1168
% MinacPasiniQuadrelliTan2020: Koszul algebras and quadratic duals in Galois cohomology; https://arxiv.org/pdf/1808.01695
% IshiiKernel2026: On the kernels of the pro-p outer Galois representations associated to once-punctured CM elliptic curves; arXiv2312.04196v3
% IshiiKummer2024: On Kummer characters arising from Galois actions on the pro-p fundamental groups of once-punctured CM elliptic curves; arXiv2312.04175v2
% McCallumSharifi2003: A cup product in the Galois cohomology of number fields; https://www.math.ucla.edu/~sharifi/publ/pairing.pdf
% BuhlerHarvey2011: Irregular primes to 163 million; arXiv0912.2121
% Iwasawa1981: Riemann-Hurwitz formula and p-adic Galois representations for number fields, Tohoku33(1981)263-288; https://www.jstage.jst.go.jp/article/tmj1949/33/2/33_2_263/_pdf/-char/en
% SharifiIwasawaNotes: Iwasawa Theory, Chapter3; https://www.math.ucla.edu/~sharifi/notes/iwasawa-ch03.html
% DeligneGoncharov2005: Groupes fondamentaux motiviques de Tate mixte; attached LaTeX Sec1ccc/Sec2ccc.
% Hirose2026: Confluence relations, Hoffman basis over Z_(2), and the Deligne-Drinfeld conjecture; arXiv2610.23719v1.
